%
% architecture.tex
%
% Lux Technical Architecture Specification
% Written for principal engineers evaluating the system.
%
% Compile: pdflatex architecture.tex
%

\documentclass[a4paper,10pt]{article}
\usepackage[margin=1in]{geometry}
\usepackage{booktabs}
\usepackage{listings}
\usepackage{xcolor}
\usepackage{enumitem}
\usepackage{longtable}
\usepackage[colorlinks=true,linkcolor=blue,citecolor=blue,urlcolor=blue]{hyperref}

\lstset{
  basicstyle=\ttfamily\small,
  keywordstyle=\color{blue!80!black},
  commentstyle=\color{green!50!black},
  stringstyle=\color{red!60!black},
  breaklines=true,
  frame=single,
  framerule=0.3pt,
  rulecolor=\color{gray!50},
  backgroundcolor=\color{gray!5},
  columns=fullflexible,
}

\newcommand{\modpath}[1]{\texttt{src/punt\_lux/#1}}
\newcommand{\tool}[1]{\texttt{#1}}
\newcommand{\msg}[1]{\texttt{#1}}
\newcommand{\elem}[1]{\texttt{#1}}
\newcommand{\field}[1]{\texttt{#1}}

\begin{document}

\title{Lux: Technical Architecture Specification}
\author{Visual Output Surface for AI Agents}
\date{May 2026 --- v0.18.0}
\hypersetup{
  pdftitle={Lux: Technical Architecture Specification},
  pdfauthor={Punt Labs},
  pdfsubject={Architecture},
  pdfcreator={pdflatex}
}
\maketitle

\tableofcontents
\newpage

% ===================================================================
\section{Overview}
% ===================================================================

Lux is a visual output surface for Claude Code.  It renders element
trees pushed by AI agents via MCP tools into a native ImGui window
over Unix socket IPC.  The system has four projection surfaces from
a single codebase: Python library, CLI, MCP server, and Claude Code
plugin.

The display server owns the render loop.  Agents are clients.  All
agent--display communication is message-passing over a framed JSON
protocol on a Unix domain socket.  The display server is
single-threaded: no threads, no asyncio, no locks.  It polls all
I/O non-blocking each frame via \texttt{select()}.

\subsection{Design Principles}

\begin{enumerate}[label=\textbf{P\arabic*}.]
  \item \textbf{Composable element tree.}  Every element kind is
    small, nestable, and data-driven.  If an agent can describe it
    as JSON, Lux renders it.
  \item \textbf{Protocol is the API surface.}  The wire protocol is
    the only interface between agents and the display.  No shared
    memory, no function calls, no Python imports required.
  \item \textbf{Single-threaded simplicity.}  The display server
    uses a polling event loop with zero threads.  All state is owned
    by the main thread.  No synchronization primitives.
  \item \textbf{Incremental by default.}  Scenes can be patched by
    element ID without replacing the entire tree.  Table filters run
    at frame rate with zero MCP round-trips.
  \item \textbf{Multi-scene coexistence.}  Each \tool{show()} call
    opens a new tab.  Scenes persist until explicitly dismissed or
    cleared.  Same \field{scene\_id} replaces content in-place.
\end{enumerate}

% ===================================================================
\section{System Architecture}
% ===================================================================

\subsection{Component Map}

\begin{center}
\begin{tabular}{llp{7cm}}
\toprule
\textbf{Module} & \textbf{File} & \textbf{Responsibility} \\
\midrule
Protocol    & \modpath{protocol.py}  & Message types, element dataclasses,
              wire format (framing, serialization, deserialization) \\
Display     & \modpath{display.py}   & ImGui render loop, socket IPC,
              scene state, widget state, element renderers, event flush \\
Client      & \modpath{client.py}    & \texttt{LuxClient} context manager;
              connects, sends, receives, buffers pending messages \\
Display Client & \modpath{display\_client.py} & High-level client with
              typed methods for show, update, introspect, screenshot \\
MCP Tools   & \modpath{tools.py}     & FastMCP tools (24 tools: show,
              update, clear, ping, recv, introspection, configuration);
              auto-reconnect \\
Hub         & \modpath{hub.py}       & WebSocket session hub; multiplexes
              MCP sessions, routes scene updates to display \\
Show        & \modpath{show.py}      & CLI subcommands for pre-built
              scenes (beads board) \\
Apps/Beads  & \modpath{apps/beads.py} & Beads issue browser; loads
              \texttt{.beads/issues.jsonl}, renders filterable table \\
Paths       & \modpath{paths.py}     & Socket discovery, PID lifecycle,
              auto-spawn \\
Runtime     & \modpath{runtime.py}   & \texttt{CodeExecutor}; compiles
              and runs render functions per-frame \\
Config      & \modpath{config.py}    & \texttt{.lux/config.md} YAML
              frontmatter; display mode toggle \\
Hooks       & \modpath{hooks.py}     & Claude Code lifecycle hook handlers \\
Service     & \modpath{service.py}   & Daemon lifecycle management
              (start, stop, status) \\
Remote      & \modpath{remote.py}    & mcp-proxy configuration and
              WebSocket transport setup \\
CLI         & \modpath{\_\_main\_\_.py} & Typer CLI: display, serve, enable,
              disable, status, doctor, install, uninstall, hub-install,
              hub-uninstall, hub-status, ensure-hub, setup-proxy, show,
              ping, version \\
\bottomrule
\end{tabular}
\end{center}

\subsection{Data Flow}

Three primary data paths connect agents to the display:

\subsubsection{Write Path (Agent $\to$ Display)}

\begin{enumerate}
  \item Agent calls MCP tool (e.g., \tool{show()}).
  \item MCP server creates \texttt{LuxClient}, calls
    \texttt{client.show(scene\_id, elements)}.
  \item Client serializes \msg{SceneMessage} as JSON, frames with
    4-byte big-endian length prefix, sends over Unix socket.
  \item Display receives frame in \texttt{\_poll\_clients()},
    deserializes, stores in \texttt{\_scenes} dict.
  \item Display renders scene each frame via
    \texttt{\_render\_scene()}.
  \item Display sends \msg{AckMessage} back to client.
\end{enumerate}

\subsubsection{Event Path (User $\to$ Agent)}

\begin{enumerate}
  \item User interacts with widget (click, slider drag, text input).
  \item Display generates \msg{InteractionMessage} with
    \field{element\_id}, \field{action}, \field{value}, and
    optional \field{scene\_id}.
  \item Message buffered in \texttt{\_event\_queue}.
  \item Each frame, \texttt{\_flush\_events()} broadcasts all queued
    messages to all connected clients, then clears the queue.
  \item Agent calls \tool{recv()} to consume events.
\end{enumerate}

\subsubsection{Update Path (Incremental Patching)}

\begin{enumerate}
  \item Agent calls \tool{update(scene\_id, patches)}.
  \item Display receives \msg{UpdateMessage}, calls
    \texttt{\_apply\_update()}.
  \item \texttt{\_find\_element()} locates target by ID via
    depth-first search through the element tree.
  \item Patch applied: \field{set} updates fields in-place,
    \field{remove} deletes element and cleans up widget state.
  \item Routes to correct scene by \field{scene\_id} (not limited
    to active tab).
\end{enumerate}

% ===================================================================
\section{Wire Protocol}
% ===================================================================

\subsection{Framing}

Every message is a length-prefixed JSON frame:

\begin{lstlisting}
+------------------+----------------------------+
| 4 bytes (uint32) |  N bytes (UTF-8 JSON)      |
| big-endian       |  payload                   |
+------------------+----------------------------+
\end{lstlisting}

Maximum payload size: 16\,MiB ($2^{24}$ bytes).  The
\texttt{FrameReader} accumulates partial reads into a byte buffer
and yields complete frames.  Oversize frames trigger client
disconnection.

\subsection{Message Types}

\subsubsection{Client $\to$ Display}

\begin{center}
\begin{tabular}{lp{9cm}}
\toprule
\textbf{Type} & \textbf{Purpose} \\
\midrule
\msg{SceneMessage}  & Replace or add scene.  Fields: \field{id},
  \field{elements[]}, \field{title}, \field{layout}. \\
\msg{UpdateMessage} & Incremental patches.  Fields: \field{scene\_id},
  \field{patches[]} (each: \field{id}, \field{set}, \field{remove}). \\
\msg{ClearMessage}  & Remove all scenes and reset state. \\
\msg{PingMessage}   & Heartbeat.  Field: \field{ts}. \\
\msg{MenuMessage}   & Set agent-provided menu entries. \\
\msg{ThemeMessage}   & Apply theme by name. \\
\msg{ConnectMessage} & Client identity handshake.  Field: \field{name}. \\
\msg{RegisterMenuMessage} & Register per-client menu items.
  Field: \field{items[]}. \\
\msg{IntrospectRequest} & Request scene element tree.
  Field: \field{scene\_id}. \\
\msg{ListScenesRequest} & Request list of all active scenes. \\
\msg{ScreenshotRequest} & Request display screenshot. \\
\msg{QueryRequest} & Generic introspection query.
  Field: \field{query\_type}. \\
\bottomrule
\end{tabular}
\end{center}

\textbf{SceneMessage frame fields} (since v0.9.0): \msg{SceneMessage}
includes optional fields \field{frame\_id}, \field{frame\_title},
\field{frame\_size} (initial \texttt{[width, height]}),
\field{frame\_flags} (dict of ImGui window flags such as
\texttt{no\_resize}, \texttt{no\_collapse}, \texttt{no\_title\_bar},
\texttt{no\_background}, \texttt{no\_scrollbar}, \texttt{auto\_resize}),
and \field{frame\_layout} (\texttt{"tab"} or \texttt{"stack"}).

\subsubsection{Display $\to$ Client}

\begin{center}
\begin{tabular}{lp{9cm}}
\toprule
\textbf{Type} & \textbf{Purpose} \\
\midrule
\msg{ReadyMessage}       & Handshake on connect; protocol version,
  capabilities. \\
\msg{AckMessage}         & Acknowledges scene or update; \field{scene\_id},
  \field{ts}, optional \field{error}. \\
\msg{InteractionMessage} & User interaction; \field{element\_id},
  \field{action}, \field{value}, \field{ts}, optional
  \field{scene\_id}. \\
\msg{PongMessage}        & Ping response; round-trip latency probe. \\
\msg{IntrospectResponse} & Scene element tree response. \\
\msg{ListScenesResponse} & List of active scene summaries. \\
\msg{ScreenshotResponse} & Base64-encoded display screenshot. \\
\msg{QueryResponse}      & Generic introspection response. \\
\bottomrule
\end{tabular}
\end{center}

Unknown message types are deserialized as \msg{UnknownMessage} for
forward compatibility rather than disconnecting the client.

\subsection{Element Kinds}

24 element kinds organized by category:

\begin{center}
\begin{longtable}{llp{6cm}}
\toprule
\textbf{Category} & \textbf{Element} & \textbf{Notes} \\
\midrule
\endfirsthead
\toprule
\textbf{Category} & \textbf{Element} & \textbf{Notes} \\
\midrule
\endhead
Text        & \elem{text}              & Content + style (heading, caption, success, error, code);
              optional \field{color} hex string \\
            & \elem{markdown}          & CommonMark rendering via imgui\_md \\
\midrule
Media       & \elem{image}             & File path or base64 data \\
\midrule
Interactive & \elem{button}            & Label, action, disabled flag;
              \field{arrow} (directional) and \field{small} variants \\
            & \elem{slider}            & Min/max/value, integer mode \\
            & \elem{checkbox}          & Boolean toggle \\
            & \elem{combo}             & Dropdown selection \\
            & \elem{input\_text}       & Single-line text input \\
            & \elem{radio}             & Radio button group \\
            & \elem{input\_number}     & Numeric input with step buttons, min/max clamping,
              integer mode \\
            & \elem{color\_picker}     & RGB/RGBA hex picker;
              \field{alpha} (RGBA) and \field{picker} (full widget) modes \\
\midrule
Data        & \elem{table}             & Columns, rows, built-in filters, detail panel, pagination \\
            & \elem{plot}              & ImPlot: line, scatter, bar series \\
            & \elem{progress}          & Fraction bar with label \\
            & \elem{spinner}           & Loading indicator \\
\midrule
Lists       & \elem{selectable}        & Click-to-select item \\
            & \elem{tree}              & Recursive node hierarchy;
              \field{flat} flag for non-indented expansion \\
\midrule
Drawing     & \elem{draw}              & Canvas with draw commands \\
\midrule
Layout      & \elem{modal}             & Modal popup dialog; container with children,
              emits \texttt{"closed"} event \\
            & \elem{group}             & Row, column, or paged layout
              (\field{pages} + \field{page\_source}) \\
            & \elem{tab\_bar}          & Tabbed container \\
            & \elem{collapsing\_header} & Expandable section \\
            & \elem{window}            & Movable, resizable floating panel \\
            & \elem{separator}         & Visual divider \\
\bottomrule
\end{longtable}
\end{center}

% ===================================================================
\section{Display Server Internals}
% ===================================================================

\subsection{Frame Cycle}

The display server runs an ImGui render loop via \texttt{immapp.run()}.
Each frame executes four phases in sequence:

\begin{enumerate}
  \item \textbf{Accept connections.}  \texttt{select()} on server
    socket with timeout~0.  Accept if readable.  Send
    \msg{ReadyMessage} to new client.
  \item \textbf{Poll clients.}  \texttt{select()} on all client
    sockets with timeout~0.  Read available data into per-client
    \texttt{FrameReader}.  Dispatch complete messages to
    \texttt{\_handle\_message()}.  Remove errored sockets.
  \item \textbf{Render scene.}  Enable docking via
    \texttt{dock\_space\_over\_viewport}.  Render idle flame background.
    Check for World menu background click.  Render World panel if open.
    Render all workspace frames (see \S\ref{sec:frames}).  Render dock
    bar for minimized frames.  Legacy fallback: if no frames, render
    unframed scenes as tabs.
  \item \textbf{Flush events.}  Broadcast all buffered
    \msg{InteractionMessage}s to all connected clients, then clear
    the queue.
\end{enumerate}

All four phases are non-blocking.  The frame completes in
microseconds when idle.  ImGui handles frame pacing via GLFW vsync.

\subsection{Multi-Scene State}

\begin{center}
\begin{tabular}{lp{8cm}}
\toprule
\textbf{Field} & \textbf{Type and Purpose} \\
\midrule
\texttt{\_scenes}              & \texttt{dict[str, SceneMessage]} --- ordered
  by insertion; keyed by \field{scene\_id}. \\
\texttt{\_scene\_order}        & \texttt{list[str]} --- explicit tab ordering. \\
\texttt{\_active\_tab}         & \texttt{str | None} --- currently selected tab. \\
\texttt{\_scene\_widget\_state} & \texttt{dict[str, WidgetState]} --- per-scene
  slider/checkbox/combo/text state. \\
\texttt{\_event\_queue}        & \texttt{list[InteractionMessage]} --- global,
  shared across all scenes. \\
\texttt{\_dirty\_windows}      & \texttt{set[str]} --- window element IDs
  needing initial position set. \\
\midrule
\multicolumn{2}{l}{\textbf{Frame state (since v0.9.0)}} \\
\midrule
\texttt{\_frames}             & \texttt{dict[str, \_Frame]} --- all workspace
  frames, keyed by \field{frame\_id}. \\
\texttt{\_focus\_frame\_id}   & \texttt{str | None} --- frame to auto-focus
  on next render. \\
\texttt{\_scene\_to\_frame}   & \texttt{dict[str, str]} --- scene
  $\to$ frame lookup. \\
\texttt{\_scene\_to\_owner}   & \texttt{dict[str, int]} --- scene
  $\to$ client fd ownership. \\
\midrule
\multicolumn{2}{l}{\textbf{Client identity (since v0.9.0)}} \\
\midrule
\texttt{\_client\_names}      & \texttt{dict[int, str]} --- client fd
  $\to$ display name. \\
\texttt{\_menu\_registrations} & \texttt{dict[int, list]} --- per-client
  menu items. \\
\texttt{\_menu\_owners}       & \texttt{dict[str, int]} --- menu item ID
  $\to$ owner fd. \\
\bottomrule
\end{tabular}
\end{center}

\subsubsection{Scene Lifecycle}

\begin{description}[style=nextline]
  \item[New scene (\field{scene\_id} not in \texttt{\_scenes})]
    Appended to \texttt{\_scenes} and \texttt{\_scene\_order}.
    Fresh \texttt{WidgetState} allocated.
    Becomes the active tab.  Window elements marked dirty.
    Existing events preserved.

  \item[Replace-in-place (same \field{scene\_id})]
    Content replaced in \texttt{\_scenes}.  Widget state and render
    function state cleared for that scene.  Events for elements
    that existed in the old scene but not the new scene are drained
    (using full subtree ID collection via \texttt{\_collect\_ids()}).
    Events from other scenes are untouched.

  \item[Dismiss (tab close button)]
    Scene removed from all state dicts.  Window element IDs removed
    from \texttt{\_dirty\_windows}.  Active tab moves to the adjacent
    neighbor (next tab, or previous if rightmost was dismissed).

  \item[Clear (\msg{ClearMessage} or menu)]
    All scenes, widget state, events,
    and dirty windows cleared.  Fresh instances allocated
    (not \texttt{.clear()} on aliased objects).
\end{description}

\subsubsection{Widget State Swap Pattern}

The 20+ element renderers all read from \texttt{self.\_widget\_state}.
Rather than threading a \field{scene\_id} parameter through every
method, the display swaps the per-scene state dict into this instance
variable before rendering each tab:

\begin{lstlisting}[language=Python]
def _render_scene_tab(self, scene_id: str) -> None:
    self._widget_state = self._scene_widget_state[scene_id]
    scene = self._scenes[scene_id]
    for elem in scene.elements:
        self._render_element(elem)
\end{lstlisting}

This is the same pattern ImGui itself uses: ``current context.''
Zero changes required in any renderer method.

\subsection{Table Filtering}

Tables support built-in filtering at frame rate with zero MCP
round-trips.  The display server renders filter controls (search
input, combo dropdowns) above the table and applies filters in the
render loop:

\begin{enumerate}
  \item Filter state stored in \texttt{WidgetState} with prefixed
    keys (e.g., \texttt{\_\_tbl\_search\_\{eid\}},\\
    \texttt{\_\_tbl\_filter\_\{col\}\_\{eid\}}).
  \item \texttt{\_apply\_table\_filters()} renders controls, snapshots
    state, detects changes.
  \item \texttt{\_filter\_indexed\_rows()} applies AND logic across
    all active filters (case-insensitive substring for search,
    exact match for combos).
  \item Pagination resets to page~0 when filters change.
  \item Detail panel renders a two-column metadata grid plus
    long-form body text for the selected row.
\end{enumerate}

% ===================================================================
\section{Frame Architecture}\label{sec:frames}
% ===================================================================

Since v0.9.0, scenes can target named \emph{frames}---persistent ImGui
windows that outlive client connections.  Frames are the primary unit
of workspace organization.

\subsection{\_Frame Dataclass}

\begin{center}
\begin{tabular}{lp{8cm}}
\toprule
\textbf{Field} & \textbf{Type and Purpose} \\
\midrule
\field{frame\_id}       & \texttt{str} --- Unique identifier. \\
\field{title}           & \texttt{str} --- Window title bar text. \\
\field{owner\_fds}      & \texttt{set[int]} --- File descriptors of contributing
  clients. \\
\field{scenes}          & \texttt{dict[str, SceneMessage]} --- All scenes in
  this frame, keyed by \field{scene\_id}. \\
\field{scene\_order}    & \texttt{list[str]} --- Tab ordering. \\
\field{active\_tab}     & \texttt{str | None} --- Currently visible scene. \\
\field{minimized}       & \texttt{bool} --- Docked to bottom bar. \\
\field{cascade\_index}  & \texttt{int} --- Stagger offset for initial
  positioning. \\
\field{initial\_size}   & \texttt{tuple[int,int] | None} --- Width/height hint
  (first use only). \\
\field{flags}           & \texttt{dict[str,bool] | None} --- ImGui window flags
  (\texttt{no\_resize}, \texttt{no\_collapse}, \texttt{no\_title\_bar},
  \texttt{no\_background}, \texttt{no\_scrollbar}, \texttt{auto\_resize}). \\
\field{layout}          & \texttt{"tab" | "stack"} --- Multi-scene composition
  mode. \\
\bottomrule
\end{tabular}
\end{center}

\subsection{Frame Lifecycle}

\begin{description}[style=nextline]
  \item[Creation]
    A frame is created on the first \msg{SceneMessage} with a
    \field{frame\_id} not already in \texttt{\_frames}.  The client's
    fd is added to \field{owner\_fds}.

  \item[Shared ownership]
    Multiple clients can contribute scenes to the same frame.  Each
    client that sends a scene targeting the frame is added to
    \field{owner\_fds}.

  \item[Orphan model]
    When a client disconnects, its scenes persist with an
    \texttt{\_ORPHAN\_FD} sentinel in place of the real fd.  The
    frame remains visible and interactive.  A new client connecting
    with the same \field{frame\_id} adopts the orphaned scenes.

  \item[Dismissal]
    The close button removes the frame and all its scenes from state.
    Widget state for all scenes in the frame is cleaned up.

  \item[Focus management]
    New frames are positioned using cascade staggering
    (\field{cascade\_index} increments per frame).  Initial size is
    applied on first render only.  \texttt{\_focus\_frame\_id} triggers
    \texttt{SetNextWindowFocus} on the next render cycle.
\end{description}

\subsection{World Menu}

Right-clicking the docking background opens the World menu panel,
providing frame management operations: list all frames, restore
minimized frames, close all frames, and access display settings.

\subsection{Dock Bar}

Minimized frames appear as clickable buttons in a bottom dock bar.
Clicking a dock bar button restores the frame to its previous
position and size.

% ===================================================================
\section{IPC and Process Management}
% ===================================================================

\subsection{Socket Discovery}

Socket path resolution, in priority order:

\begin{enumerate}
  \item \texttt{\$LUX\_SOCKET} environment variable.
  \item \texttt{\$XDG\_RUNTIME\_DIR/lux/display.sock}.
  \item \texttt{/tmp/lux-\$USER/display.sock} (fallback).
\end{enumerate}

The fallback path is chosen to stay within macOS's $\sim$104-character
\texttt{AF\_UNIX} path limit.

\subsection{PID Lifecycle}

\begin{itemize}
  \item PID file written at \texttt{\{socket\_path\}.pid} on server
    startup.
  \item Clients verify liveness via \texttt{os.kill(pid, 0)} before
    connecting.
  \item Stale socket cleanup: \texttt{cleanup\_stale\_socket()} removes
    socket and PID file if the process is dead.
\end{itemize}

\subsection{Auto-Spawn}

\texttt{ensure\_display()} spawns the display server if not running:

\begin{enumerate}
  \item Checks socket existence and PID liveness.
  \item If not running, spawns\\
    \texttt{python -m punt\_lux display -{}-socket \{path\}}\\
    with \texttt{start\_new\_session=True} (detached from caller's
    process group).
  \item Polls socket existence + PID liveness every 100\,ms until
    ready or 5\,s timeout.
  \item Display logs to \texttt{\{socket\_path\}.log}.
\end{enumerate}

\subsection{MCP Server Reconnect}

The MCP server wraps every tool call in
\texttt{\_with\_reconnect()}: if an \texttt{OSError} (broken pipe)
occurs, the stale client is closed, a new connection is established,
and the call is retried exactly once.  No retry loop.

% ===================================================================
\section{Performance}
% ===================================================================

\subsection{Frame Budget}

\begin{center}
\begin{tabular}{lrl}
\toprule
\textbf{Parameter} & \textbf{Value} & \textbf{Notes} \\
\midrule
Idle FPS            & 30             & \texttt{fps\_idling.fps\_idle = 30.0} \\
Active FPS          & 60+            & GLFW vsync, GPU-bound \\
Socket poll timeout & 0              & Non-blocking \texttt{select()} \\
Max message size    & 16\,MiB        & Hard limit in \texttt{FrameReader} \\
Connection timeout  & 5.0\,s         & \texttt{ensure\_display()} and handshake \\
Texture loading     & Lazy           & First render triggers PIL + OpenGL upload \\
\bottomrule
\end{tabular}
\end{center}

\subsection{Critical Path Analysis}

The frame-critical path is:

\begin{enumerate}
  \item \texttt{select()} on $N$ client sockets --- $O(N)$, $N < 64$
    in practice.
  \item \texttt{FrameReader.feed()} --- memcpy into buffer, yield
    frames.  $O(B)$ where $B$ is bytes received.
  \item \texttt{\_render\_scene()} --- ImGui draw list calls.
    $O(E)$ where $E$ is total element count across visible scene.
  \item \texttt{\_flush\_events()} --- send $K$ events to $N$
    clients.  $O(K \cdot N)$.
\end{enumerate}

Table filtering is $O(R \cdot C)$ where $R$ is row count and $C$
is column count.  This runs every frame but is fast for tables
under $\sim$10{,}000 rows.  No benchmarks exist for larger tables;
pagination (configurable page size) bounds visible row rendering.

\subsection{Idle Screen}

The idle screen renders $\sim$60 draw calls per frame: 48 radial
lines, 3 glow circles, 3 flame shapes (each $\sim$7 bezier
operations).  This is negligible for a GPU-rendered immediate-mode
UI at 30\,FPS idle.

\subsection{Memory}

\begin{itemize}
  \item \textbf{Scene storage.}  Scenes are Python dataclass trees
    held in memory.  No size limit beyond the 16\,MiB message cap.
    Each scene is a separate dict entry; dismissed scenes are
    garbage-collected.
  \item \textbf{Texture cache.}  OpenGL texture IDs held until
    \texttt{\_on\_exit()}.  No eviction policy.  Image-heavy scenes
    accumulate GPU memory.
  \item \textbf{Widget state.}  Per-scene \texttt{WidgetState} is a
    flat dict.  Cleared on scene replacement.  No growth concern.
  \item \textbf{Event queue.}  Flushed every frame.  Bounded by
    interaction rate ($<$100 events/frame in practice).
\end{itemize}

\subsection{Known Limitations}

\begin{enumerate}
  \item \textbf{No ImGui ID scoping.}  Widget IDs are derived from
    \field{element\_id} alone, not scoped by scene.  ImGui internal
    state (focus, open/closed) can bleed between tabs with colliding
    IDs.
  \item \textbf{No texture eviction.}  Long-running sessions with
    many images accumulate GPU memory.
  \item \textbf{Filter is per-frame.}  Table filtering re-executes
    every frame even when nothing changed.  Not a problem at current
    scale but would benefit from dirty-flag optimization for
    10{,}000+ row tables.
\end{enumerate}

% ===================================================================
\section{Security}
% ===================================================================

\subsection{Threat Model}

The primary trust boundary is the Unix socket.  The display server
accepts connections from any process that can reach the socket file.
In the default configuration (\texttt{/tmp/lux-\$USER/}), the socket
is protected by filesystem permissions (user-only directory).

\begin{center}
\begin{tabular}{lp{5cm}p{5cm}}
\toprule
\textbf{Threat} & \textbf{Description} & \textbf{Mitigation} \\
\midrule
Socket hijacking    & Another process connects to the socket and
  injects scenes or reads events. & Directory permissions (0700) on
  \texttt{/tmp/lux-\$USER/}. \\
Malformed messages  & Client sends invalid JSON, unknown message type,
  or oversize payload. & Unknown types deserialized as
  \msg{UnknownMessage}; oversize payloads rejected by FrameReader. \\
Render function     & Agent-submitted Python code executes arbitrary
  operations. & Consent dialog with AST warnings.  User must
  explicitly click Allow. \\
Resource exhaustion & Scene with millions of elements or very large
  images. & 16\,MiB message cap.  No element count limit (ImGui
  handles gracefully). \\
\bottomrule
\end{tabular}
\end{center}

\subsection{Render Function Security}

Render functions are the only element kind that executes arbitrary
code.  The security model is \textbf{consent-based, not
sandbox-based}:

\begin{enumerate}
  \item \textbf{AST scanning}: a best-effort visitor flags
    suspicious imports (\texttt{os}, \texttt{subprocess},
    \texttt{socket}, \texttt{ctypes}, \texttt{urllib}, etc.) and
    dangerous builtins (\texttt{eval}, \texttt{exec}, \texttt{open},
    \texttt{\_\_import\_\_}).  This is a \emph{UX signal}, not a
    security boundary.

  \item \textbf{Consent dialog}: an ImGui modal displays the full
    source code with line numbers and any AST warnings.  The user
    must click ``Allow'' to proceed.  Dismissing the modal (Esc) is
    treated as ``Deny.''  Consent is per-element: changing the source
    code resets consent.

  \item \textbf{Execution}: \texttt{CodeExecutor} compiles the
    source and calls \texttt{render(ctx)} each frame.  The
    \texttt{RenderContext} provides a persistent state dict, delta
    time, frame counter, canvas dimensions, and a \texttt{send()}
    callback for emitting events.  Execution is \textbf{not sandboxed}:
    the code runs in the display server's process with full access to
    the Python runtime.  Runtime errors are caught and displayed, not
    propagated.
\end{enumerate}

\textbf{Implication}: render functions should be treated as
user-approved code, equivalent to running a script the user
downloaded.  The consent dialog provides informed consent, not
isolation.

\subsection{Protocol Robustness}

\begin{itemize}
  \item \textbf{Malformed JSON}: client disconnected, not crashed.
    The display logs the error and removes the client from
    \texttt{\_clients}.
  \item \textbf{Unknown message type}: deserialized as
    \msg{UnknownMessage} for forward compatibility.
  \item \textbf{Oversize frame}: \texttt{FrameReader} checks
    \texttt{bufSize $\leq$ maxBufSize} and disconnects on violation.
  \item \textbf{Partial reads}: \texttt{FrameReader} accumulates
    bytes across frames; no assumption of complete messages per
    \texttt{recv()}.
  \item \textbf{Identity fields protected}: \texttt{\_apply\_update()}
    silently drops patches that attempt to change \field{id} or
    \field{kind} on an element.
  \item \textbf{Double-remove idempotent}:
    \texttt{\_remove\_client()} is safe to call twice on the same
    socket.
\end{itemize}

% ===================================================================
\section{Platform Integration}
% ===================================================================

\subsection{macOS}

\begin{itemize}
  \item \textbf{Dock hiding.}  After GLFW initialization
    (in \texttt{\_on\_post\_init()}), the display server sets\\
    \texttt{NSApplicationActivationPolicyAccessory}\\
    on the shared \texttt{NSApplication}.  GLFW's
    \texttt{glfwInit()} sets the policy to Regular; the post-init
    call overrides it.  The Dock icon briefly flashes at launch.
    Requires \texttt{pyobjc-framework-Cocoa} (optional
    \texttt{display} extra).
  \item \textbf{Process name.}  \texttt{setproctitle("Lux")} makes
    the process identifiable in \texttt{ps} and \texttt{top}.  Does
    not affect Activity Monitor's binary name column.
  \item \textbf{Font discovery.}  Prefers Arial Unicode or Helvetica;
    merges Apple Symbols for mathematical/Z notation glyphs.
  \item \textbf{AF\_UNIX path limit.}  macOS limits Unix socket
    paths to $\sim$104 characters.\\
    \texttt{tempfile.mkdtemp(prefix="lux-")}
    used in tests; production paths are short by design.
\end{itemize}

\subsection{Linux}

\begin{itemize}
  \item \textbf{Font discovery.}  DejaVu Sans or Noto Sans; merges
    Noto Sans Symbols2 for Unicode coverage.
  \item \textbf{XDG compliance.}  Prefers
    \texttt{\$XDG\_RUNTIME\_DIR/lux/} for socket placement.
  \item \textbf{No Dock equivalent.}  Window managers use the X11
    window title (set to ``Lux'').
\end{itemize}

% ===================================================================
\section{MCP Tool Surface}
% ===================================================================

24 tools organized by function:

\begin{center}
\begin{longtable}{lp{8cm}}
\toprule
\textbf{Tool} & \textbf{Signature and Purpose} \\
\midrule
\endfirsthead
\toprule
\textbf{Tool} & \textbf{Signature and Purpose} \\
\midrule
\endhead
\multicolumn{2}{l}{\textbf{Scene management}} \\
\midrule
\tool{show}           & \texttt{(scene\_id, elements, title?, layout?,
  frame\_id?, frame\_title?, frame\_size?, frame\_flags?)}
  --- Display a scene.  Docstring lists all 24 element kinds. \\
\tool{show\_table}    & \texttt{(scene\_id, columns, rows, filters?,
  detail?, title?)} --- Filterable data table with search, combo
  filters, and detail panel. \\
\tool{show\_dashboard} & \texttt{(scene\_id, metrics?, charts?,
  table\_columns?, table\_rows?, title?)} --- Metrics cards + charts +
  summary table. \\
\tool{update}         & \texttt{(scene\_id, patches)} --- Incremental
  JSON patches targeting elements by ID. \\
\tool{clear}          & \texttt{()} --- Clear all scenes. \\
\midrule
\multicolumn{2}{l}{\textbf{Communication}} \\
\midrule
\tool{ping}           & \texttt{()} --- Latency probe; returns
  \msg{PongMessage}. \\
\tool{recv}           & \texttt{(timeout?)} --- Receive next event
  (interaction, window lifecycle). \\
\tool{set\_menu}      & \texttt{(menus)} --- Set agent-provided menu
  bar entries. \\
\tool{register\_tool} & \texttt{(name, description, schema)} --- Register
  a custom tool in the display's tool palette. \\
\tool{set\_theme}     & \texttt{(theme)} --- Apply display theme. \\
\midrule
\multicolumn{2}{l}{\textbf{Configuration}} \\
\midrule
\tool{display\_mode}  & \texttt{()} --- Read current display mode
  (y/n); advisory L3 signal for consumer plugins. \\
\tool{set\_display\_mode} & \texttt{(mode)} --- Set display mode. \\
\tool{set\_window\_settings} & \texttt{(opacity?, font\_scale?,
  decorated?, fps\_idle?)} --- Configure display window appearance
  and idle frame rate. \\
\tool{set\_frame\_state} & \texttt{(frame\_id, minimized?)}
  --- Control frame state (minimize/restore). \\
\midrule
\multicolumn{2}{l}{\textbf{Introspection}} \\
\midrule
\tool{inspect\_scene} & \texttt{(scene\_id)} --- Return element tree
  for a scene. \\
\tool{list\_scenes}   & \texttt{()} --- List all active scenes with
  metadata. \\
\tool{screenshot}     & \texttt{()} --- Capture display as base64
  PNG. \\
\tool{get\_display\_info} & \texttt{()} --- Display dimensions,
  frame count, client count. \\
\tool{get\_window\_settings} & \texttt{()} --- Current window
  geometry. \\
\tool{get\_theme}     & \texttt{()} --- Current theme name. \\
\tool{list\_clients}  & \texttt{()} --- Connected clients with
  names and scene counts. \\
\tool{list\_menus}    & \texttt{()} --- Registered menu items. \\
\tool{list\_recent\_events} & \texttt{(count?)} --- Recent
  interaction events. \\
\tool{list\_errors}   & \texttt{(count?)} --- Recent error log
  entries. \\
\bottomrule
\end{longtable}
\end{center}

% ===================================================================
\section{Formal Model}
% ===================================================================

The display server's state machine is formally specified in
\texttt{docs/display-server.tex} using the Z notation (ISO 13568).
The specification covers:

\begin{itemize}
  \item Client connection/disconnection with FrameReader association.
  \item Scene replacement, incremental patching, and clearing.
  \item Interaction event generation and broadcast.
  \item FrameReader streaming buffer invariants.
\end{itemize}

The specification has been type-checked with Spivey's \texttt{fuzz}
and model-checked with ProB (7{,}943 reachable states explored).
Refinement tests in \texttt{tests/test\_display\_refinement.py}
verify commutativity between the abstract Z operations and the
concrete Python implementation via an abstraction function
(\texttt{tests/display\_abstraction.py}).

Partition tests in \texttt{tests/test\_display\_partition.py} derive
test cases from the Z specification using the Test Template Framework
(TTF), covering all input space partitions of each operation schema.

\textbf{Note}: the formal model predates multi-scene tabs and models
a single-scene display.  The abstraction function maps multi-scene
state to the active tab for backward compatibility.  A specification
update to model the full multi-scene state is pending.

% ===================================================================
\section{Test Architecture}
% ===================================================================

\begin{center}
\begin{tabular}{lrl}
\toprule
\textbf{Test File} & \textbf{Count} & \textbf{What It Tests} \\
\midrule
\texttt{test\_protocol.py}            & 168 & Serialization, deserialization, framing \\
\texttt{test\_display\_partition.py}   & 107 & TTF-derived partitions from Z spec \\
\texttt{test\_tools.py}               & 83  & MCP tool behavior (mocked client) \\
\texttt{test\_display\_state.py}       & 65  & State machine logic (no ImGui) \\
\texttt{test\_display\_refinement.py}  & 28  & Abstraction function commutativity \\
\texttt{test\_display\_client.py}      & 26  & High-level client library \\
\texttt{test\_query.py}               & 23  & Generic introspection queries \\
\texttt{test\_show.py}                & 22  & CLI show subcommands \\
\texttt{test\_runtime.py}             & 22  & CodeExecutor lifecycle \\
\texttt{test\_hooks.py}               & 18  & Hook handler output \\
\texttt{test\_config.py}              & 16  & Config read/write \\
\texttt{test\_paths.py}               & 16  & Socket discovery, PID, auto-spawn \\
\texttt{test\_remote.py}              & 12  & mcp-proxy and remote transport \\
\texttt{test\_cli.py}                 & 11  & CLI command smoke tests \\
\texttt{test\_paths\_hub.py}          & 11  & Hub socket and PID paths \\
\texttt{test\_service.py}             & 9   & Daemon lifecycle \\
\texttt{test\_unit.py}                & 6   & Misc unit tests \\
\texttt{test\_hub.py}                 & 6   & WebSocket session hub \\
\texttt{test\_version.py}             & 2   & Version string consistency \\
\midrule
\multicolumn{3}{l}{\textbf{Marker-gated (not in default \texttt{pytest} run)}} \\
\midrule
\texttt{test\_integration.py}         & 8   & Socket IPC, protocol framing (real sockets) \\
\texttt{test\_e2e.py}                 & 3   & Full display server lifecycle (requires GPU) \\
\midrule
\textbf{Total}                        & 662 & (651 default + 11 marker-gated) \\
\bottomrule
\end{tabular}
\end{center}

Quality gates: \texttt{ruff} (lint + format), \texttt{mypy} (strict),
\texttt{pyright} (strict), \texttt{pytest} (all markers except
slow/integration/e2e).  Zero violations required.

\end{document}
